% Generated draft: requires mathematical and source review.
\documentclass[11pt]{article}
\usepackage[margin=25mm]{geometry}
\usepackage{fontspec}
\setmainfont{Latin Modern Roman}
\usepackage{amsmath,amssymb,mathtools,unicode-math}
\setmathfont{Latin Modern Math}
\usepackage{enumitem,needspace,xurl,hyperref}
\setlength{\parindent}{0pt}
\setlength{\parskip}{5pt}
\setlength{\emergencystretch}{3em}
\newcommand{\N}{\mathbb{N}}
\newcommand{\Z}{\mathbb{Z}}
\newcommand{\Q}{\mathbb{Q}}
\newcommand{\R}{\mathbb{R}}
\newcommand{\C}{\mathbb{C}}
\newcommand{\F}{\mathbb{F}}
\newcommand{\A}{\mathbb{A}}
\newcommand{\PP}{\mathbb{P}}
\newcommand{\E}{\mathbb{E}}
\newcommand{\eps}{\varepsilon}
\newcommand{\dd}{\,\mathrm{d}}
\begin{document}
\section*{UnsolvedMath ID: 11000201}
\textbf{Source problem code:} catalogue entry
\subsection{Definitions and mathematical statement}
Let \(\mathcal T(S)=\mathcal T_{g,m}\) be the intended Teichm\"uller space, \(\pi:\mathcal T(S)\to\mathcal M(S)\) the moduli projection, and \(\Psi:\mathcal T(S)\to\mathcal C(S)\) the curve-complex map defined in the source's preceding discussion. Write \(d\) for the target metric. The excerpt does not supply the choice or coarse conventions of \(\Psi\), the surface complexity or the parameterization of geodesics. In the customary unit-Teichm\"uller-speed interpretation, let \(\gamma:\R\to\mathcal T(S)\) be a complete geodesic and define the exact source condition
\[d\bigl(\Psi(\gamma(s)),\Psi(\gamma(t))\bigr)\geq\frac{|s-t|}{R}-R\qquad(s,t\in\R).\tag{*}\]
Only a lower bound is stated; no additional upper inequality is inserted.

The compound question asks:
\begin{enumerate}
\item For every fixed \(R>0\), is there a compact \(K(R)\subset\mathcal M(S)\), independent of \(\gamma\), such that every geodesic satisfying (*) has \(\pi(\gamma(\R))\subseteq K(R)\)?
\item Conversely, for every compact \(K\subset\mathcal M(S)\), is there \(R(K)>0\), independent of the geodesic, such that every complete geodesic with \(\pi(\gamma(\R))\subseteq K\) satisfies (*) for all real \(s,t\)? The source's notation ``Mod(\(S\))'' in this clause is interpreted as moduli space in light of its compact-set wording, not as a new sampling measure on the mapping class group.
\item Analyze the \(\mathcal C(S)\)-images of Teichm\"uller geodesics associated with minimal filling geodesic laminations that are not uniquely ergodic. Minimal means every leaf is dense in the support; filling means the support meets every essential curve; nonunique ergodicity permits more than one nonproportional transverse measure on that support. The measure, transverse data and basepoint needed to determine an actual geodesic or ray must be specified: an unmeasured support alone does not single out one parameterized geodesic. No precise dynamical classification of these images is supplied.
\end{enumerate}
The two compactness/coarse-progress implications have different uniformity quantifiers, both retained. The source does not say that a compact set may vary with each geodesic in the first clause, or that \(R\) may vary with each geodesic in the second.
\subsection{Short English statement}
Ask whether uniform linear progress in the curve complex is equivalent, with the stated uniform constants, to staying in a compact part of moduli space. Also study the curve-complex images associated with filling but nonuniquely ergodic laminations.
\subsection{Sources}
\begin{itemize}
\item Local UnsolvedMath registry, permanent ID \(11000201\), exact source code \texttt{catalogue entry}.
\item Farb, \emph{Problems on Mapping Class Groups and Related Topics}, source heading ``Problem 2:'', PDF page 242: \url{https://www.math.uchicago.edu/~farb/papers/mcgbook.pdf}. Registry access 2026-07-29; pdf-font-structure extraction. No separate definition or result citation is supplied.
\end{itemize}
\end{document}
